% !PRINTED 567
\weberpage{567}

%<<CONT
woraus:
\weq{7}{a\partial - bc = \dfrac{x^2 - \Delta y^2}{4} = \mu\mu' = m,}
wenn
\weq{8}{\mu' = \dfrac{x - y\sqrt{\Delta}}{2}}
die zu $\mu$ konjugierte Zahl des Körpers $\Omega$ ist.

Nehmen wir $y$ von Null verschieden, so folgt aus (1) durch Multiplikation mit $y$:
\[ b\omega_2^2 + (a-\partial)\omega_2\omega_1 - c\omega_1^2 = 0, \]
und dafür kann man auch schreiben:
\[ (a\omega_1 + b\omega_2)\omega_2 = (c\omega_1 + \partial\omega_2)\omega_1 \]
oder
\weq{9}{\omega = \dfrac{c + \partial\omega}{a + b\omega}.}
Nach (4) und (5) folgt aus (6):
\weq{10}{\begin{aligned} a + b\omega &= \dfrac{x + (2A\omega+B)y}{2} = \mu, \\ c + \partial\omega &= \mu\omega, \end{aligned}}
oder in homogener Form:
\weq{11}{\begin{aligned} \mu\omega_1 &= a\omega_1 + b\omega_2, \\ \mu\omega_2 &= c\omega_1 + \partial\omega_2, \end{aligned}}
und durch Auflösung dieser linearen Gleichungen:
\weq{12}{\begin{aligned} \mu'\omega_1 &= \phantom{-}\partial\omega_1 - b\omega_2, \\ \mu'\omega_2 &= -c\omega_1 + a\omega_2. \end{aligned}}

Daraus ergibt sich, daß $\wp(\mu u)$ eine doppelt-periodische Funktion mit den Perioden $\omega_1$, $\omega_2$ ist, und da sie außerdem eine gerade Funktion von $u$ ist, so läßt sie sich rational durch $\wp(u)$ darstellen (§ 21). Wenn wir also mit $R$ und $P$ ganze rationale Funktionen von $\wp(u)$ ohne gemeinschaftlichen Teiler bezeichnen, so ist
\weq{13}{\wp(\mu u) = \dfrac{R}{P}.}

Da der Quotient $\wp(\mu u):\wp(u)$ für $u=0$, d.\,h. für $\wp(u)=\infty$, endlich bleibt, so ist der Grad von $R$ um eine Einheit höher als der Grad des Nenners.
